\begin{abstract}
The SAE J1962 diagnostic connector is increasingly used as the permanent
mount and power source of telematics, insurance and fleet devices, yet it
was designed for occasional workshop use. Vibration-induced fretting
corrosion of its contacts causes intermittent data loss and device
brown-outs that are routinely misattributed to the device or the vehicle.
This paper asks whether connector wear can be detected, attributed and
predicted from the telemetry the attached tool already records. We propose
excitation-coherent telemetry analysis (ECTA): self-referenced loss and
loss-structure features plus a Poisson test of whether losses co-vary with
the vehicle's own measured engine and road excitation, as a fretting
contact's interruptions must. A physics-informed model maps contact wear to
the observable behaviour of full-bus loggers, sampling loggers and OBD-II
polling testers, including four realistic confounders. All healthy
telemetry, excitation and usage are real, from four public datasets
covering five vehicles in detail and a 381-vehicle fleet (236,817 labelled
windows), evaluated with vehicle-disjoint splits against ten baseline and ablation models.
Severe wear is detected at every tool fidelity (AUROC 0.85--0.98 against
healthy telemetry; 0.999 per vehicle after pooling ten trips), moderate
wear only with full-bus capture (0.875). Excitation coherence is what
separates wear from bus interference in polling mode (0.947 vs. 0.737) and
raises fleet-wide detection from 0.704 to 0.775. A learned health index
yields early warnings for 88\% of failing connectors with 1.5\% false
alarms, but not months-ahead remaining-life estimates for low-rate tools.
\end{abstract}

\begin{keywords}
Controller area network, electrical contacts, fault diagnosis, fretting
corrosion, intelligent vehicles, on-board diagnostics, predictive
maintenance, prognostics and health management, remaining useful life,
vehicle telematics.
\end{keywords}
